%% Drafted from the mapped corpus by scripts/draft_topics.py.
% Model: gpt-5.6-terra; source characters: 29744.
\section{Параметрични уравнения на права}

Нека в афинната равнина е фиксирана афинна координатна система
$K=Oxy$. Права $g$ се определя еднозначно от точка $M_0$ от нея и
ненулев вектор $\vec p$, колинеарен на правата. Векторът $\vec p$ се
нарича направляващ вектор на $g$.

\begin{defn}
Нека $M_0\in g$ и $\vec p\neq\vec 0$, $\vec p\parallel g$. Векторното
параметрично уравнение на правата $g$ е
\[
 \vec r=\vec r_0+s\vec p,\qquad s\in\mathbb R,
\]
където $\vec r=\overrightarrow{OM}$ е радиус-векторът на текущата точка
$M$, а $\vec r_0=\overrightarrow{OM_0}$ е радиус-векторът на фиксираната
точка $M_0$.
\end{defn}

Действително,
\[
 M\in g
 \quad\Longleftrightarrow\quad
 \overrightarrow{M_0M}\parallel\vec p
 \quad\Longleftrightarrow\quad
 \exists s\in\mathbb R:\ \overrightarrow{M_0M}=s\vec p.
\]
Понеже
\[
 \overrightarrow{M_0M}
 =\overrightarrow{OM}-\overrightarrow{OM_0}
 =\vec r-\vec r_0,
\]
получаваме векторното параметрично уравнение.

\begin{defn}
Нека $M_0(x_0,y_0)$ и $\vec p=(p_1,p_2)\neq(0,0)$. Координатните
(скаларни) параметрични уравнения на правата през $M_0$, колинеарна на
$\vec p$, са
\[
 g:\quad
 \begin{cases}
 x=x_0+sp_1,\\
 y=y_0+sp_2,
 \end{cases}
 \qquad s\in\mathbb R.
\]
\end{defn}

Параметърът $s$ задава положението на точката върху правата. При $s=0$
се получава началната точка $M_0$. Различни параметризации могат да
задават една и съща права: може да се избере друга точка от правата или
ненулев колинеарен на $\vec p$ вектор.

\begin{example}
Правата през $M_0(2,-1)$ с направляващ вектор
$\vec p=(3,2)$ има уравнения
\[
 \vec r=(2,-1)+s(3,2),\qquad s\in\mathbb R,
\]
съответно
\[
 \begin{cases}
 x=2+3s,\\
 y=-1+2s,
 \end{cases}
 \qquad s\in\mathbb R.
\]
Например при $s=1$ се получава точката $(5,1)$.
\end{example}

\section{Общо уравнение на права}

\begin{keythm}[Теорема за общото уравнение на права]
Всяка права в равнината има уравнение от вида
\[
 Ax+By+C=0,\qquad (A,B)\neq(0,0).
\]
Обратно, всяко уравнение от този вид определя точно една права в
равнината.
\end{keythm}

\begin{proof}
Нека правата $g$ минава през $M_0(x_0,y_0)$ и има направляващ вектор
$\vec p=(p_1,p_2)\neq(0,0)$. Точка $M(x,y)$ лежи върху $g$ тогава и само
тогава, когато векторите
\[
 \overrightarrow{M_0M}=(x-x_0,y-y_0)
 \quad\text{и}\quad
 \vec p=(p_1,p_2)
\]
са колинеарни. Координатното условие за колинеарност е
\[
 \begin{vmatrix}
 x-x_0 & y-y_0\\
 p_1 & p_2
 \end{vmatrix}=0.
\]
Следователно
\[
 p_2(x-x_0)-p_1(y-y_0)=0,
\]
тоест
\[
 p_2x-p_1y-p_2x_0+p_1y_0=0.
\]
При полагането
\[
 A=p_2,\qquad B=-p_1,\qquad C=-p_2x_0+p_1y_0
\]
получаваме уравнение от търсения вид. Понеже $\vec p\neq\vec0$, то
$(A,B)=(p_2,-p_1)\neq(0,0)$.

Обратно, нека е дадено
\[
 Ax+By+C=0,\qquad (A,B)\neq(0,0).
\]
Нека $M_0(x_0,y_0)$ е негова точка. Разглеждаме вектора
\[
 \vec p=(-B,A)\neq\vec0.
\]
Съществува единствена права през $M_0$, колинеарна на $\vec p$. Нейните
параметрични уравнения са
\[
 x=x_0-sB,\qquad y=y_0+sA,\qquad s\in\mathbb R.
\]
Оттук
\[
 A(x-x_0)+B(y-y_0)=0.
\]
Тъй като $Ax_0+By_0+C=0$, последното уравнение е еквивалентно на
$Ax+By+C=0$. Следователно даденото уравнение задава точно тази права.
\end{proof}

\begin{defn}
Уравнението
\[
 g:\ Ax+By+C=0,\qquad (A,B)\neq(0,0),
\]
се нарича общо уравнение на правата $g$ спрямо координатната система
$K=Oxy$.
\end{defn}

Коефициентите на общото уравнение дават два важни вектора:
\[
 \vec n=(A,B)\perp g,\qquad
 \vec p=(-B,A)\parallel g.
\]
Векторът $\vec n$ се нарича нормален вектор на правата, а
$\vec p$ е един от нейните направляващи вектори.

\begin{prop}
Нека
\[
 g:\ Ax+By+C=0,\qquad (A,B)\neq(0,0).
\]
Векторът $\vec q=(q_1,q_2)$ е колинеарен на $g$ тогава и само тогава,
когато
\[
 Aq_1+Bq_2=0.
\]
\end{prop}

\begin{proof}
Направляващ вектор на $g$ е $\vec p=(-B,A)$. Условието
$\vec q\parallel\vec p$ е
\[
 \begin{vmatrix}
 -B & A\\
 q_1 & q_2
 \end{vmatrix}=0,
\]
което е равносилно на $Aq_1+Bq_2=0$.
\end{proof}

\begin{example}
Да се намери общото уравнение на правата през $M_0(1,2)$ с
направляващ вектор $\vec p=(4,-3)$.

От доказателството на теоремата получаваме
\[
 A=p_2=-3,\qquad B=-p_1=-4,
\]
а
\[
 C=-p_2x_0+p_1y_0=3+8=11.
\]
Следователно
\[
 -3x-4y+11=0,
\]
или еквивалентно
\[
 3x+4y-11=0.
\]
\end{example}

\section{Взаимно положение на две прави}

Нека са дадени правите
\[
 g_1:\ A_1x+B_1y+C_1=0,\qquad
 g_2:\ A_2x+B_2y+C_2=0,
\]
където
\[
 (A_i,B_i)\neq(0,0),\qquad i=1,2.
\]
Техните общи точки са решенията на системата
\[
 \begin{cases}
 A_1x+B_1y+C_1=0,\\
 A_2x+B_2y+C_2=0.
 \end{cases}
\]
Определителят на системата по неизвестните е
\[
 \Delta=
 \begin{vmatrix}
 A_1&B_1\\
 A_2&B_2
 \end{vmatrix}
 =A_1B_2-A_2B_1.
\]

\begin{keythm}
За взаимното положение на $g_1$ и $g_2$ са в сила следните критерии.

\begin{enumerate}
 \item Правите се пресичат в точно една точка тогава и само тогава,
 когато
 \[
 \Delta\neq0.
 \]

 \item Правите съвпадат тогава и само тогава, когато съществува
 $k\neq0$, за което
 \[
 A_1=kA_2,\qquad B_1=kB_2,\qquad C_1=kC_2.
 \]

 \item Правите са различни успоредни тогава и само тогава, когато
 съществува $k\neq0$, за което
 \[
 A_1=kA_2,\qquad B_1=kB_2,\qquad C_1\neq kC_2.
 \]
\end{enumerate}
\end{keythm}

\begin{proof}
Ако $\Delta\neq0$, системата има единствено решение, следователно
правите имат точно една обща точка.

Ако $\Delta=0$, двойките $(A_1,B_1)$ и $(A_2,B_2)$ са пропорционални.
Тогава или и свободните членове са пропорционални със същия коефициент,
което прави двете уравнения еквивалентни, или не са. В първия случай
правите съвпадат, а във втория системата е несъвместима и правите са
различни успоредни.
\end{proof}

\begin{remark}
Записите с отношения, например
\[
 \frac{A_1}{A_2}=\frac{B_1}{B_2},
\]
са удобни само когато съответните знаменатели са различни от нула.
Критериите чрез общ ненулев множител $k$ или чрез определителя
$\Delta$ нямат това ограничение.
\end{remark}

\begin{supp}[Бележка]
Среща се неточен критерий, според който пресичането се описва с
неравенства между всички отношения на коефициентите. Точният критерий
за пресичане е
\[
 A_1B_2-A_2B_1\neq0.
\]
Свободните членове $C_1,C_2$ определят положението на пресечната точка,
но не влияят на това дали две прави с различни направления се
пресичат.
\end{supp}

\begin{example}
За правите
\[
 g_1:\ 2x-y+1=0,\qquad
 g_2:\ 4x-2y-3=0
\]
е изпълнено
\[
 \Delta=
 \begin{vmatrix}
 2&-1\\
 4&-2
 \end{vmatrix}=0.
\]
Първите два коефициента на $g_2$ са два пъти съответните коефициенти на
$g_1$, но свободният член $-3$ не е два пъти $1$. Следователно правите
са различни успоредни.
\end{example}

\section{Декартово уравнение}

Оттук нататък при метричните понятия приемаме, че $K=Oxy$ е
ортонормирана координатна система.

\begin{defn}
Нека
\[
 g:\ Ax+By+C=0,\qquad B\neq0.
\]
След изразяване на $y$ получаваме
\[
 y=-\frac ABx-\frac CB.
\]
Ако
\[
 k=-\frac AB,\qquad n=-\frac CB,
\]
то уравнението
\[
 g:\ y=kx+n
\]
се нарича декартово уравнение на правата $g$.
\end{defn}

Коефициентът $k$ се нарича ъглов коефициент на правата. Ако
ориентираният ъгъл между положителната посока на $Ox$ и правата е
$\alpha$, то
\[
 k=\tan\alpha.
\]
Един направляващ вектор на правата $y=kx+n$ е $(1,k)$.

\begin{prop}
Права има декартово уравнение $y=kx+n$ тогава и само тогава, когато не
е успоредна на координатната ос $Oy$.
\end{prop}

\begin{proof}
Декартовото уравнение се получава точно когато $B\neq0$. Ако $B=0$,
общото уравнение има вид $Ax+C=0$, тоест описва права, успоредна на
$Oy$. Обратно, когато $B\neq0$, $y$ се изразява еднозначно чрез $x$.
\end{proof}

\begin{remark}
Правата $y=kx+n$ пресича оста $Oy$ в точката $(0,n)$. Тя пресича
оста $Ox$, ако $k\neq0$, в точката
\[
 \left(-\frac nk,0\right).
\]
\end{remark}

\begin{supp}[Бележка]
В някои материали точката $(0,n)$ е означена като пресечна точка с
$Ox$. Това е несъответствие: при уравнение $y=kx+n$ тази точка лежи на
оста $Oy$.
\end{supp}

\begin{example}
Правата
\[
 3x-2y+6=0
\]
има декартово уравнение
\[
 y=\frac32x+3.
\]
Тя има ъглов коефициент $k=\frac32$ и пресича $Oy$ в точката $(0,3)$.
\end{example}

\section{Нормално уравнение и разстояние от точка до права}

Нека
\[
 g:\ Ax+By+C=0,\qquad (A,B)\neq(0,0).
\]
Нормален вектор на правата е $\vec n=(A,B)$, а неговата дължина е
\[
 |\vec n|=\sqrt{A^2+B^2}.
\]
Следователно един от двата единични нормални вектора е
\[
 \vec n_1=
 \left(
 \frac{A}{\sqrt{A^2+B^2}},
 \frac{B}{\sqrt{A^2+B^2}}
 \right).
\]

\begin{defn}
Нормално уравнение на правата $g$ се нарича общо уравнение на същата
права, при което коефициентите пред $x$ и $y$ са координати на единичен
нормален вектор. Нормалните уравнения на $g$ са
\[
 \frac{Ax+By+C}{\sqrt{A^2+B^2}}=0
 \qquad\text{и}\qquad
 -\frac{Ax+By+C}{\sqrt{A^2+B^2}}=0.
\]
\end{defn}

Следователно всяка права има точно две нормални уравнения. Те се
получават чрез двата противоположно насочени единични нормални вектора.

\begin{keythm}[Разстояние от точка до права]
Нека
\[
 g:\ Ax+By+C=0,\qquad (A,B)\neq(0,0),
\]
и $P(x_1,y_1)$ е точка в равнината. Тогава разстоянието от $P$ до $g$
е
\[
 d(P,g)=
 \frac{|Ax_1+By_1+C|}{\sqrt{A^2+B^2}}.
\]
\end{keythm}

\begin{proof}
Нека нормалното уравнение на $g$ е
\[
 A_1x+B_1y+C_1=0,
\]
където $A_1^2+B_1^2=1$. Нека $H(x_H,y_H)$ е ортогоналната проекция на
$P$ върху $g$. Тогава $\overrightarrow{HP}$ е колинеарен на единичния
нормален вектор $(A_1,B_1)$, тоест за точно едно $\delta\in\mathbb R$
е изпълнено
\[
 \overrightarrow{HP}=\delta(A_1,B_1).
\]
Следователно
\[
 x_H=x_1-\delta A_1,\qquad
 y_H=y_1-\delta B_1.
\]
Тъй като $H\in g$, имаме
\[
 A_1x_H+B_1y_H+C_1=0.
\]
След заместване и използване на $A_1^2+B_1^2=1$ получаваме
\[
 \delta=A_1x_1+B_1y_1+C_1.
\]
Числото $\delta$ е ориентираното разстояние, а търсеното разстояние е
неговата абсолютна стойност. След заместване на $A_1,B_1,C_1$ се
получава посочената формула.
\end{proof}

\begin{remark}
Ако уравнението на правата вече е нормално, тоест
\[
 A^2+B^2=1,
\]
формулата се опростява до
\[
 d(P,g)=|Ax_1+By_1+C|.
\]
Без абсолютна стойност изразът
\[
 \delta(P,g)=\frac{Ax_1+By_1+C}{\sqrt{A^2+B^2}}
\]
е ориентираното разстояние, определено от избрания единичен нормален
вектор.
\end{remark}

\begin{example}
Да се намери разстоянието от $P(1,-2)$ до правата
\[
 g:\ 3x+4y-5=0.
\]
Изчисляваме
\[
 d(P,g)=
 \frac{|3\cdot1+4\cdot(-2)-5|}{\sqrt{3^2+4^2}}
 =\frac{|-10|}{5}=2.
\]
\end{example}

\section{Полуравнини}

Всяка права разделя равнината на две части, наречени отворени
полуравнини. Нека
\[
 g:\ Ax+By+C=0,\qquad (A,B)\neq(0,0),
\]
и въведем линейната функция
\[
 l(x,y)=Ax+By+C.
\]
За точките от правата е изпълнено $l(x,y)=0$. За всяка точка извън
правата стойността на $l(x,y)$ е или положителна, или отрицателна.

\begin{keythm}
Нека $M_1(x_1,y_1)$ и $M_2(x_2,y_2)$ не лежат на правата $g$. Те лежат
в различни полуравнини спрямо $g$ тогава и само тогава, когато
\[
 l(x_1,y_1)\,l(x_2,y_2)<0.
\]
\end{keythm}

\begin{proof}
Нека първо $M_1$ и $M_2$ са в различни полуравнини. Тогава отсечката
$M_1M_2$ пресича $g$ в точка $M_0(x_0,y_0),$ която е вътрешна за
отсечката. Следователно съществува $k<0$, за което
\[
 \overrightarrow{M_0M_1}=k\overrightarrow{M_0M_2}.
\]
Оттук
\[
 x_0=\frac{x_1-kx_2}{1-k},\qquad
 y_0=\frac{y_1-ky_2}{1-k}.
\]
Тъй като $M_0\in g$, след заместване в $l(x_0,y_0)=0$ получаваме
\[
 l(x_1,y_1)-k\,l(x_2,y_2)=0.
\]
Следователно
\[
 k=\frac{l(x_1,y_1)}{l(x_2,y_2)}<0,
\]
откъдето
\[
 l(x_1,y_1)l(x_2,y_2)<0.
\]

Обратно, нека
\[
 l(x_1,y_1)l(x_2,y_2)<0.
\]
Тогава
\[
 k=\frac{l(x_1,y_1)}{l(x_2,y_2)}<0.
\]
Точката
\[
 M_0\left(
 \frac{x_1-kx_2}{1-k},
 \frac{y_1-ky_2}{1-k}
 \right)
\]
удовлетворява $l(x_0,y_0)=0$, следователно лежи на $g$. Освен това
\[
 \overrightarrow{M_0M_1}
 =k\overrightarrow{M_0M_2},\qquad k<0.
\]
Векторите са противопосочни, значи $M_0$ е вътрешна точка на
отсечката $M_1M_2$. Следователно $M_1$ и $M_2$ са в различни
полуравнини.
\end{proof}

\begin{defn}
Двете отворени полуравнини с граница $g$ се задават аналитично чрез
\[
 \lambda_g=\{M(x,y):Ax+By+C>0\}
\]
и
\[
 \overline{\lambda}_g=\{M(x,y):Ax+By+C<0\}.
\]
\end{defn}

Знакът на $l(x,y)$ е постоянен във всяка от двете полуравнини. Това
дава практичен начин за проверка от коя страна на права се намира
дадена точка.

\begin{example}
Нека
\[
 g:\ 2x-y-1=0.
\]
За точките $P(0,0)$ и $Q(2,0)$ имаме
\[
 l(0,0)=-1,\qquad l(2,0)=3.
\]
Техният произведение е отрицателно, следователно $P$ и $Q$ лежат в
различни полуравнини спрямо $g$.
\end{example}

\section{Изпитен фокус}

\begin{itemize}
 \item Да се дефинира права чрез точка $M_0$ и ненулев направляващ
 вектор $\vec p$, както и да се запише векторно-параметричното
 уравнение
 \[
 \vec r=\vec r_0+s\vec p.
 \]

 \item Да се преминава към координатни параметрични уравнения
 \[
 x=x_0+sp_1,\qquad y=y_0+sp_2.
 \]

 \item Да се формулира и докаже теоремата за общото уравнение
 \[
 Ax+By+C=0,\qquad (A,B)\neq(0,0),
 \]
 в двете посоки, чрез условието за колинеарност и вектора
 $(-B,A)$.

 \item Да се разпознават нормален вектор $(A,B)$ и направляващ вектор
 $(-B,A)$ на права с общо уравнение.

 \item Да се определя взаимното положение на две прави чрез
 \[
 \Delta=A_1B_2-A_2B_1
 \]
 и чрез пропорционалността на тройките коефициенти.

 \item Да се извежда декартовото уравнение
 \[
 y=kx+n,\qquad k=-\frac AB,\qquad n=-\frac CB,
 \]
 при $B\neq0$, и да се обясни защо правите, успоредни на $Oy$, нямат
 такава форма.

 \item Да се нормализира общо уравнение на права и да се запишат двете
 нормални уравнения.

 \item Да се възпроизвежда формулата
 \[
 d(P,g)=\frac{|Ax_1+By_1+C|}{\sqrt{A^2+B^2}}
 \]
 и идеята за доказателството чрез ортогоналната проекция на точката
 върху правата.

 \item Да се описват полуравнините чрез знака на
 $l(x,y)=Ax+By+C$ и да се прилага критерият
 \[
 l(x_1,y_1)l(x_2,y_2)<0
 \]
 за точки, разположени в различни полуравнини.
\end{itemize}
